\section{Lean formalization and the verification boundary}\label{sec:formalization}

The archive contains a concrete formalization outline, including 24 candidate
closed derivations in the object language of \ipra.  \emph{The new Lean sources
were not typechecked in this delivery: no Lean or Lake executable was available
in the execution environment.}  The human proofs in this report and the
formalization plan therefore have a different status from a completed,
kernel-checked formalization.  In particular, the numerical checks do not change
the Lean status.  The files contain no \texttt{sorry}, \texttt{admit}, or new
unproved axiom declaration.

\subsection{The actual target theory}

The inspected ProveIt revision is
\texttt{bbb85238dac533fc4b7b84a750cca864fe35fed5}.  Its relevant syntax is not a
surrogate geometry theory: it is the ordinary first-order presentation of
\ipra, implemented in
\href{https://github.com/VladimirReshetnikov/ProveIt/blob/bbb85238dac533fc4b7b84a750cca864fe35fed5/Logic/PrimitiveRecursiveArithmetic/ConciseIPRA/Lean/ConciseIPRA/PR.lean}{\texttt{ConciseIPRA/PR.lean}}.
The constructors have the following meanings.

\begin{center}
\begin{tabular}{p{0.24\linewidth}p{0.66\linewidth}}
\texttt{PR k} & A specific primitive-recursive function symbol of arity $k$,
generated by zero, successor, projections, composition, and primitive recursion.\\[3pt]
\texttt{PTm n} & A term over $n$ natural-number variables, formed by applying
those symbols to terms.\\[3pt]
\texttt{Fm PTm n} & An intuitionistic first-order formula over those variables.\\[3pt]
\texttt{PRAAx} & The defining equations, the successor axioms, and induction for
quantifier-free formulas.\\[3pt]
\texttt{Prf PRAAx Gamma phi} & An explicit derivation of the object formula
\texttt{phi} from the context \texttt{Gamma}.\\
\end{tabular}
\end{center}

An endpoint in the desired theory has the form
\[
   \texttt{Prf PRAAx [] phi},\qquad
   \texttt{phi : Fm PTm 0}.
\]
Proving a theorem about Mathlib's real numbers does not by itself produce this
object.  We use Mathlib in the metatheory to construct an \ipra\ derivation;
the axioms available to that derivation remain precisely
\texttt{PRAAx}.  The repository's own
\href{https://github.com/VladimirReshetnikov/ProveIt/blob/bbb85238dac533fc4b7b84a750cca864fe35fed5/Logic/PrimitiveRecursiveArithmetic/README.md}{status document}
also distinguishes its checked derivation library from delivered research
packets that have not been rerun or verified.

\subsection{Natural-number definitions underlying the geometry}

The Lean data model is in \path{lean/ExactCodes.lean}.  It uses a redundant signed
pair instead of the even/odd integer code of Section~\ref{sec:foundations};
the two representations have explicit primitive-recursive conversions by sign
comparison and truncated subtraction.  A signed integer is a pair
$(p,n)$ representing $p-n$.  We do not demand a canonical representation.
Addition, negation, and multiplication are
\[
 (p,n)+(r,s)=(p+r,n+s),\qquad -(p,n)=(n,p),
 \qquad (p,n)(r,s)=(pr+ns,ps+nr).
\]
Equality of the represented integers is the natural-number equality
$p+s=r+n$; order is $p+s\le r+n$.  Thus the apparent use of negative coordinates
introduces no second numerical sort into \ipra.

A rational code is a triple $(p,n,d)$ denoting $(p-n)/(d+1)$.
The denominator is positive by construction.  For two such triples the
predecessor of the product denominator is
\[
 d_1d_2+d_1+d_2,
 \qquad
 (d_1d_2+d_1+d_2)+1=(d_1+1)(d_2+1).
\]
The numerator of a sum is formed by cross multiplication; multiplication uses
the signed-product rule above.  Equality and order compare the cross products
with positive denominators.  Inversion first compares $p$ and $n$.  If $p=n$,
the total program returns zero; if $p>n$, it returns the positive numerator
$d+1$ and denominator $p-n$; if $p<n$, it uses numerator $-(d+1)$ and denominator
$n-p$.  Every geometric invocation of inversion includes the relevant nonzero
guard.  The zero branch is only a way of making a total \PR\ function.

Points are records of two or three rational codes.  A line has three
coefficients, a plane has four, and a planar quadratic has six.  Finite tuples
can be encoded by iterated natural-number pairing.  Consequently these records
are notation for fixed-length natural data, and their equality, incidence,
and rational sign tests are \PR.  The file also gives the Cramer intersection
program and the square-root index
\[
  j(A,B,k)=\operatorname{isqrt}\!\left(\left\lfloor
                A4^k/B\right\rfloor\right),\qquad B>0.
\]
The output bracket is $j/2^k,(j+1)/2^k$, as in Theorem~\ref{found:sqrt}.
Its \ipra\ correctness should be
obtained from the inspected square-root module below, rather than from an
unrelated semantic theorem about a real-valued square-root function.

\subsection{Explicit analytic programs}

\path{lean/AnalyticPrograms.lean} adds the concrete finite programs needed for
the later analytic chapters.  A sparse polynomial is a list of rational
coefficients and finite exponent vectors.  Evaluation is a finite sum of
monomials; a formal partial derivative multiplies a coefficient by the chosen
exponent and decreases that exponent by one.  A zero exponent produces a zero
coefficient, so its truncated predecessor cannot introduce an extra term.
The functions \texttt{H}, \texttt{B}, and \texttt{L} implement exactly the three
coefficient bounds used in Section~\ref{sec:tangency}.

The trigonometric programs sum the even and odd terms through total degree $N$.
They implement $N=4M+n+6$ for rational input, and the stated argument-query and
degree schedules for a fixed named input.  The implicit-curve program stores
the vertical coordinate first, explicitly permuting the report's $F(x,y)$ to
the code's ordering.  It evaluates the polynomial exactly at the midpoint,
branches on its rational sign, and performs precisely $n+2$ iterations.
Its correctness still requires the root, positive-partial-derivative, and box
guards from Theorem~\ref{tan:implicit}; the program is total even when those
guards do not hold.

For a polynomial curve the file supplies the coefficient bounds $V,K$, dyadic
partition points, exact squared chord lengths, and the rational arc-length
observer.  With $C=\lceil2K(b-a)^2\rceil+1$, it uses
$k=n+\ell(C)+3$, $2^k$ chords, and square-root precision $n+k+3$ for each chord,
matching Theorem~\ref{tan:arclength}.  These definitions fill in the computational
objects behind the outline.  They do not assert that their correctness has
already been formalized.

\subsection{Twenty-four explicit polynomial endpoints}

The central new source is \path{lean/ObjectCertificates.lean}.  For an integer
term $u$, the repository uses a record \texttt{ZT n} with natural-number terms
\texttt{pos} and \texttt{neg}.  The formula \texttt{ZT.eq u v} is precisely the
cross-addition equality above.  To state an identity in $r$ arbitrary signed
integer coordinates, the file introduces $2r$ natural-number variables,
groups them into signed pairs, writes the equality as an object formula, and
universally closes all $2r$ variables with \texttt{closeUniv}.

Each proof invokes \texttt{ipra\_ring}, whose implementation was inspected in
\href{https://github.com/VladimirReshetnikov/ProveIt/blob/bbb85238dac533fc4b7b84a750cca864fe35fed5/Logic/PrimitiveRecursiveArithmetic/Derivations/Lean/IPRADerivations/Core/Tactic.lean}{\texttt{Core/Tactic.lean}}.
This tactic works through a specific derivability quotient.  First form the
quotient of natural-number terms by equality provable from the current
\ipra\ context.  The arithmetic library makes this quotient a cancellative
commutative semiring.  Passing to its ring of differences gives a commutative
ring in which signed polynomial identities can be normalized.  The theorem
\texttt{ZT.derive} then turns the normalized equality back into an object-level
\texttt{Prf PRAAx} derivation.  This last step is why the tactic is suitable:
normalization is not being mistaken for a general semantic reflection axiom.

The 24 endpoints are listed below.  Several vector statements contain two or
three coordinate identities in a single conjunction; they count as one named
endpoint.  The names and exact arities are also recorded in
\path{lean/certificate_index.json}.

\begin{center}
\begin{tabular}{p{0.36\linewidth}p{0.56\linewidth}}
\texttt{lagrange2}, \texttt{parallelogram}, \texttt{cosine}
 & Gram determinant, parallelogram law, and squared cosine law.\\[4pt]
\texttt{cramer\_row1}, \texttt{cramer\_row2}, \texttt{cramer\_recover}
 & Cramer numerators solve the scaled equations; the adjugate recovers both
 coordinates.\\[4pt]
\texttt{rotation\_norm}, \texttt{rotation\_dot}, \texttt{affine\_determinant}
 & Norm and dot-product transport; oriented area scaling.\\[4pt]
\texttt{reflection\_norm}, \texttt{projection\_remainder}
 & Cleared-denominator reflection and projection identities.\\[4pt]
\texttt{circle\_parameter}, \texttt{sphere\_parameter}
 & Rational circle and stereographic sphere parametrizations.\\[4pt]
\texttt{quadratic\_completion}, \texttt{parabola\_focus}
 & Completing the square and the focus--directrix equation.\\[4pt]
\texttt{conic\_line\_expansion}, \texttt{discriminant\_congruence}
 & Exact quadratic restriction to a line; discriminant transport.\\[4pt]
\texttt{cross\_orthogonal}, \texttt{lagrange3}, \texttt{vector\_triple}
 & Orthogonality, the three-dimensional Gram identity, and the vector triple
 product.\\[4pt]
\texttt{skew\_line\_invariant}, \texttt{weighted\_variance}
 & Normal separation of skew lines; the three-point variance identity.\\[4pt]
\texttt{radical\_axis}, \texttt{conic\_midpoint}
 & Difference of circle powers; symmetric second differences of quadratics.\\
\end{tabular}
\end{center}

For example, the reflection endpoint states an unconditional polynomial
identity.  With $q=a^2+b^2$ and $h=ax+by$ it is
\[
 (qx-2ha)^2+(qy-2hb)^2=q^2(x^2+y^2).
\]
The geometric reflection theorem divides by $q^2$ only after proving $q>0$.
Keeping the unconditional identity and the guarded division separate makes
the remaining proof obligation visible and permits reuse at degenerate input.
Similarly, a Cramer identity is not already a theorem asserting the existence
of a unique intersection: existence needs the constructed quotient, and
uniqueness needs cancellation under the determinant guard.

\subsection{Reusable inspected derivations}

The following are concrete links to existing source, not claims that the whole
new report was imported or rerun.  Their paths and Git blob hashes are recorded
in \path{lean/consulted_sources.json}.

\begin{enumerate}
\item \texttt{Core/Integers.lean} supplies signed terms and
\texttt{ZT.derive}; \texttt{Core/Tactic.lean} supplies the polynomial tactic and
universal closure.
\item \path{Geometry/PlaneMetric.lean} already contains
\texttt{gramClosed} and \texttt{midpointClosed}.  These can replace the
corresponding fresh candidate certificates after the signatures are matched.
\item \texttt{Geometry/TriangleFormulas.lean} contains closed derivations of
the two-square identity, Heron's coordinate identity, and the cosine law.
\item \texttt{Geometry/Circumcircle.lean} derives the scaled circumcenter
equidistance identities and its radius identity.  It supplies an existing
algebraic base for the circumcenter construction in the report.
\item \texttt{Geometry/PowerOfPoint.lean} contains signed secant, chord, and
tangent--secant derivations.  Their guards should be preserved when importing
them into a new coordinate formulation.
\item The dyadic square-root module~\cite{repo-sqrt} contains
\texttt{index\_bracket}, \texttt{index\_greatest}, and closed sentences for the
actual dyadic bracket, source bound, next endpoint, and output program.  In
particular its bracket is about the explicit integer-square-root computation,
with the positive-denominator hypothesis.
\end{enumerate}

Reusing a packet under \texttt{docs/proofs} would require a separate audit of its
proofs.  The references above instead identify concrete derivation source in
\texttt{Derivations/Lean}; even these imported sources must be built in the
chosen checkout before a local verification claim is made.

\subsection{The rational and represented-real bridge}

\path{lean/BridgeSpec.lean} defines a symbolic rational term as an integer
numerator and a natural denominator.  Its validity formula is $1\le d$.
Rational equality is the integer equation $ad'=a'd$.  The file supplies
symbolic addition and multiplication and a structure listing their required
validity and congruence derivations.  That structure has no supplied instance;
its fields are a work list, not assumptions silently granted to the geometry.

The bridge is completed in the following order.

\begin{enumerate}
\item Prove that the triple coding in \texttt{ExactCodes} agrees with the
chosen native rational representation.  Prove positivity of every constructed
denominator, invariance under equivalent codes, exact cross-product comparison,
and the guarded inverse equations.  All algorithms are fixed \PR\ operations.
\item For each rational polynomial assertion, take the product of all input
denominators, raised to a sufficiently large finite degree, as a common positive
multiplier.  The resulting signed polynomial identity is one of the preceding
endpoints, or is produced by the same tactic.  Cancel the multiplier using its
positive guard.  For inequalities use the corresponding integer-order library.
\item For classifications, compile rational comparisons into finite Boolean
branches.  Carry an affine map and its inverse as output data, and prove the
normal-form equation in every branch.  Merely returning a label such as
``ellipse'' is not the required theorem.  Empty, point, repeated-line, and
two-line branches must remain present.
\item For lengths and algebraic intersection coordinates, connect the explicit
square-root bracket to the numerator/denominator representation.  Apply the
report's error estimates at its stated input precision.  A reciprocal consumes
an explicit apartness witness; square root at zero instead uses its quadratic
continuity modulus.
\item For coordinate maps and tangent estimates, instantiate the fixed
polynomial or fixed transcendental \PR\ approximants and carry the finite
remainder estimates.  The parameters, degree, coefficient bounds, denominator
gaps, and requested precision are all natural-number inputs of the final
program and formula.
\end{enumerate}

A represented-real family in the sense of Definition~\ref{found:names} is
formalized by three \emph{specific} \PR\ symbols
for the signed numerator and denominator at each precision.  The
\texttt{FixedFamily} record in \texttt{BridgeSpec} has exactly this meaning.
It does not contain a universal evaluator for arbitrary \PR\ program codes.
The family binder is part of the Lean metatheory; one first chooses the three
programs and then emits an ordinary first-order sentence about their
natural-number parameters.  This distinction avoids both second-order
quantification and the false claim that a universal total \PR\ evaluator is
primitive recursive.

The data file's \texttt{NameNonnegative} and \texttt{GapContract} match the
report's contracts.  Its finite apartness witness is a sufficient input format:
if $|X(k+2)|\ge2^{-k}+2^{-k-2}$, then the Cauchy condition and the rational
reverse triangle inequality give
\[
 |X(n)|\ge |X(k+2)|-|X(n)-X(k+2)|
       \ge2^{-k}-2^{-n}
\]
at every $n$.  Thus a finite witness supplies the report's all-precision gap
without deciding whether an arbitrary name is zero.

For each theorem and each requested precision $k$, the proof specifies a finite
input demand and a \PR\ output.  Write $H(\vec a,k)$ for the finite conjunction
of needed rational observation conditions, and $C(\vec a,k)$ for the desired
output inequality.  The final task is to construct
\[
 \ipra\vdash\forall\vec a\,\forall k\,
       \bigl(H(\vec a,k)\longrightarrow C(\vec a,k)\bigr).
\]
Proposition~\ref{found:transfer} supplies the finite-expression translation used
here.  Each matrix is expressed by equations and order tests of \PR\ terms; bounded
finite lists of tests are compressed using their characteristic functions.
If a theorem has a universal name-coherence hypothesis, its proof extracts only
the finitely many instances needed at that demand.  Induction is applied to the
quantifier-free invariant of the finite computation, not to an arbitrary
formula quantifying over all future precisions.  This is the precise bridge
between the report's readable name notation and the induction restriction.

\subsection{Build and acceptance gates}

The target checkout uses the following pinned providers:
\begin{itemize}
\item Lean \texttt{leanprover/lean4:v4.32.0}.
\item Mathlib commit \texttt{81a5d257c8e410db227a6665ed08f64fea08e997}.
\end{itemize}
The archive does not duplicate
the repository and its transitive dependencies.  With those dependencies
installed, run \path{lean/check_in_proveit.sh} with the absolute path of the
pinned ProveIt checkout.  It verifies the commit, builds the dependency
libraries, invokes Lean separately on the four new files, and records the
printed axiom dependencies of all 24 closed candidates.

Acceptance has three distinct gates.  First, each source must elaborate and
typecheck, and its axiom report must contain no unexpected assumption such as
\texttt{sorryAx}.  Second, the rational and analytic interfaces must receive
actual proofs with the stated bounds; an uninstantiated obligation record is
not sufficient.  Third, the completed closed endpoint must be entered in the
repository's axiom audit with its exact statement.  The supplied script itself
passes a shell syntax check, but the Lean build gates remain unexecuted.
These explicit remaining tasks give a direct route to a full formalization
without inflating the verification status of the current packet.
